\documentclass[11pt,a4paper]{article}

\usepackage[T1]{fontenc}
\usepackage[utf8]{inputenc}
\usepackage[french]{babel}
\usepackage[margin=2.5cm]{geometry}
\usepackage{booktabs}
\usepackage{siunitx}
\usepackage[table]{xcolor}
\usepackage{hyperref}
\usepackage{pgfplots}
\pgfplotsset{compat=1.18}
\usetikzlibrary{positioning}
\usepackage{caption}
\usepackage{enumitem}
\usepackage{amsmath}
\usepackage{amssymb}
\usepackage{listings}

\sisetup{detect-weight=true, detect-family=true}

\lstset{
  basicstyle=\ttfamily\small,
  breaklines=true,
  frame=single,
  columns=fullflexible,
  backgroundcolor=\color{black!3}
}

\hypersetup{
  colorlinks=true,
  linkcolor=blue!50!black,
  urlcolor=blue!50!black,
  citecolor=blue!50!black,
  pdftitle={Scan H(P) Al4C3 - preparation},
  pdfauthor={Pr. John Akapulco}
}

\title{\vspace{-1.5cm}Rapport de préparation\\[2pt]
\large Balayage enthalpie--pression H(P) des polymorphes Al$_4$C$_3$ (0--100~GPa)}
\author{Pr. John Akapulco \\ \small \href{mailto:computchem86@gmail.com}{computchem86@gmail.com}}
\date{22 août 2026}

\begin{document}
\maketitle
\vspace{-0.5cm}

\begin{abstract}
\noindent
Ce rapport documente la préparation d'une campagne de calculs DFT (VASP, \texttt{ISIF=8})
visant à tracer l'enthalpie $H(P) = E + PV$ des sept polymorphes d'Al$_4$C$_3$ déjà relaxés à
50~GPa, sur la gamme \textbf{0--100~GPa} (pas de 10~GPa), afin de déterminer le domaine de
stabilité de chaque phase sur toute cette gamme. Chaque polymorphe donne lieu à une
\textbf{chaîne} de 10 calculs partant d'un point « graine » à 50~GPa (ré-optimisé à
\texttt{KSPACING=0.2} avec sauvegarde de la fonction d'onde et de la densité de charge), puis
redescendant jusqu'à la borne \textbf{0~GPa} et remontant jusqu'à la borne \textbf{100~GPa} par
pas de 10~GPa, chaque point redémarrant électroniquement (\texttt{WAVECAR}), en densité de
charge (\texttt{CHGCAR}) et ioniquement (\texttt{CONTCAR}) depuis le point de pression voisin
déjà calculé dans la même branche. Soit \textbf{77 calculs} au total (7 phases $\times$ 11
pressions : 0, 10, 20, ..., 100~GPa), organisés en \textbf{7 jobs SLURM} — un par phase —
réservant chacun \textbf{un seul nœud pendant toute la durée de la série} : les 11 points de
la chaîne s'exécutent séquentiellement dans cette même allocation, le nœud n'étant libéré
qu'une fois les 11 points de la phase terminés. Les scripts sont générés et prêts dans
\texttt{hp\_curves/} ; \textbf{aucun calcul n'a encore été soumis} — la soumission attend la
validation de l'utilisateur avant de lancer la campagne sur le cluster.
\end{abstract}

\section{Objectif}

La campagne précédente (voir \texttt{reports/al4c3\_activity/}) a établi un renversement de
stabilité entre 0 et 50~GPa parmi sept polymorphes candidats d'Al$_4$C$_3$ : R-3m est le plus
stable à 0~GPa mais devient le moins stable des sept à 50~GPa, au profit de Pbam et Pnma. Pour
localiser ce croisement et cartographier le domaine de stabilité de chaque phase sur toute la
gamme 0--100~GPa, cette nouvelle campagne calcule l'enthalpie $H(P)$ de chaque polymorphe à 10
pressions supplémentaires (0, 10, 20, 30, 40, 60, 70, 80, 90, 100~GPa, pas de 10~GPa), en plus
du point 50~GPa déjà disponible qui sert de graine à la chaîne.

\section{Structures de départ}

Chaque chaîne part de la structure relaxée à 50~GPa lors de la campagne précédente
(\texttt{KSPACING=0.03}) :

\begin{table}[h]
\centering
\begin{tabular}{llc}
\toprule
Polymorphe & Source (CONTCAR.relaxed) & Atomes/maille \\
\midrule
I-43d            & \texttt{structures/Al4C3\_I-43d\_50GPa/}      & 28 \\
Pnma-III         & \texttt{structures/Al4C3\_Pnma-III\_50GPa/}   & 28 \\
Pnma             & \texttt{structures/Al4C3\_Pnma\_50GPa/}       & 28 \\
Pnma-VII         & \texttt{structures/Al4C3\_Pnma-VII\_50GPa/}   & 28 \\
R-3m             & \texttt{structures/Al4C3\_R-3m\_50GPa/}       & 21 \\
Pbam             & \texttt{structures/Al4C3\_Pbam/}              & 14 \\
anti-CaFe$_2$O$_4$ & \texttt{structures/Al4C3\_anti-CaFe2O4/}    & 28 \\
\bottomrule
\end{tabular}
\caption{Les sept polymorphes relaxés à 50~GPa (campagne précédente) servant de point de
départ à ce balayage.}
\end{table}

\section{Méthodologie}

\subsection{Schéma de chaînage}

Pourquoi une chaîne plutôt que 77 calculs indépendants partant chacun de la structure 50~GPa
d'origine ? Parce que l'énoncé demande explicitement de réutiliser, à chaque pression, la
fonction d'onde (\texttt{WAVECAR}), la densité de charge (\texttt{CHGCAR}) et la structure
(\texttt{CONTCAR}) optimisées à la pression voisine déjà calculée, pour accélérer la
convergence électronique à chaque nouveau point. Or \textbf{aucun \texttt{WAVECAR}/\texttt{CHGCAR}
n'existe dans les relaxations 50~GPa déjà réalisées} (\texttt{LWAVE=.FALSE.} et
\texttt{LCHARG=.FALSE.} dans tous les \texttt{INCAR.relax} d'origine, fichiers à 0 octet), et
un \texttt{WAVECAR}/\texttt{CHGCAR} calculé à \texttt{KSPACING=0.03} ne serait de toute façon
pas réutilisable à \texttt{KSPACING=0.2} (maillage \textbf{k} différent). Chaque phase démarre
donc par un point « graine » à 50~GPa, ré-optimisé à \texttt{KSPACING=0.2} avec
\texttt{LWAVE=.TRUE.} et \texttt{LCHARG=.TRUE.}, à partir duquel les deux branches
(descendante jusqu'à 0~GPa, ascendante jusqu'à 100~GPa) sont chaînées :

\begin{center}
\begin{tikzpicture}[node distance=13mm, every node/.style={font=\small}]
  \node (p0)  [draw, thin, circle, inner sep=1pt] {0};
  \node (p10) [right=of p0] {10};
  \node (p20) [right=of p10] {20};
  \node (p30) [right=of p20] {30};
  \node (p40) [right=of p30] {40};
  \node (p50) [right=of p40, draw, thick, circle, inner sep=1pt] {50};
  \node (p60) [right=of p50] {60};
  \node (p70) [right=of p60] {70};
  \node (p80) [right=of p70] {80};
  \node (p90) [right=of p80] {90};
  \node (p100)[right=of p90, draw, thin, circle, inner sep=1pt] {100};
  \draw[-latex] (p50) -- (p40);
  \draw[-latex] (p40) -- (p30);
  \draw[-latex] (p30) -- (p20);
  \draw[-latex] (p20) -- (p10);
  \draw[-latex] (p10) -- (p0);
  \draw[-latex] (p50) -- (p60);
  \draw[-latex] (p60) -- (p70);
  \draw[-latex] (p70) -- (p80);
  \draw[-latex] (p80) -- (p90);
  \draw[-latex] (p90) -- (p100);
\end{tikzpicture}
\end{center}

Chaque flèche représente une étape séquentielle \textbf{à l'intérieur d'un même job SLURM} :
contrairement à une première version de ce script qui soumettait un job indépendant par point
de pression (chaînés entre eux par \texttt{--dependency=afterok}, sans garantie de tomber sur
le même nœud physique d'un point à l'autre), \textbf{un seul job réserve un unique nœud pour
toute la série d'une phase}, et exécute les 11 points les uns après les autres dans cette même
allocation — le point suivant hérite directement de \texttt{WAVECAR}, \texttt{CHGCAR} et
\texttt{CONTCAR} sur le disque local à l'étape qui vient de se terminer, sans repasser par la
file d'attente SLURM entre deux points. Le point 50~GPa (graine, cerclé en gras) est commun aux
deux branches ; les bornes 0 et 100~GPa (cerclées en fin de chaîne) sont les extrémités de la
gamme demandée. Cela donne, par phase, 1 graine + 5 points descendants (40, 30, 20, 10, 0) + 5
points ascendants (60, 70, 80, 90, 100) = 11 calculs exécutés en séquence sur un seul nœud,
soit $7 \times 11 = \mathbf{77}$ calculs (sur \textbf{7 nœuds réservés en parallèle}, un par
phase) pour l'ensemble de la campagne.

\subsection{Paramètres DFT}

Les paramètres électroniques/ioniques communs reprennent ceux de la campagne précédente
(PBE, \texttt{ENCUT=600}, \texttt{ISIF=8}, \texttt{EDIFFG=-0.01}), à l'exception du maillage
\textbf{k}, volontairement relâché pour cette campagne à 77 calculs :

\begin{center}
\begin{tabular}{ll}
\toprule
Paramètre & Valeur \\
\midrule
\texttt{KSPACING / KGAMMA} & \textbf{0.2}~\AA$^{-1}$ / .TRUE. \\
\texttt{IBRION / ISIF} & 2 / 8 \\
\texttt{NSW / EDIFFG} & 200 / $-0.01$~eV/\AA \\
\texttt{PSTRESS} & 0 à 1000~kBar (0 à 100~GPa), pas de 100~kBar \\
\texttt{KPAR} & \textbf{8} (pas de \texttt{NPAR}, demandé explicitement) \\
\texttt{LWAVE / LCHARG} & \textbf{.TRUE.} / \textbf{.TRUE.} \\
\texttt{ISTART} & 0 (graine 50~GPa) ; 1 (tous les autres points) \\
\texttt{ICHARG} & non défini (graine, défaut 2) ; \textbf{1} (points chaînés, lit \texttt{CHGCAR}) \\
\bottomrule
\end{tabular}
\end{center}

\texttt{INCAR.relax} du point graine (exemple, R-3m à 50~GPa) :
\lstinputlisting{../../hp_curves/Al4C3_R-3m/P50/INCAR.relax}

\texttt{INCAR.relax} d'un point chaîné (exemple, R-3m à 40~GPa, hérite de P50) :
\lstinputlisting{../../hp_curves/Al4C3_R-3m/P40/INCAR.relax}

\subsection{Scripts de soumission SLURM}

Chaque phase possède un unique script \texttt{hp\_curves/Al4C3\_<phase>/run\_chain.sh} : c'est
le mécanisme de soumission principal. Il réserve \textbf{48 cœurs sur un seul nœud}
(\texttt{--ntasks=48}, ce qui cible naturellement les nœuds \texttt{node01--node08}, seule
famille du cluster disposant de 48~CPU logiques) avec une limite de temps de 72~h (la partition
\texttt{defq} n'a pas de limite de temps maximale), puis exécute en interne, l'un après
l'autre, les 11 points de pression de la phase — chaque étape copiant
\texttt{CONTCAR}/\texttt{WAVECAR}/\texttt{CHGCAR} du répertoire voisin juste terminé avant de
relancer \texttt{vasp\_std} sur le même nœud :

\texttt{run\_chain.sh} (exemple, R-3m — un seul job pour les 11 points de la phase) :
\lstinputlisting{../../hp_curves/Al4C3_R-3m/run_chain.sh}

Chaque répertoire \texttt{hp\_curves/Al4C3\_<phase>/P<pression>/} conserve en plus son propre
\texttt{job.sh} individuel, \textbf{gardé uniquement comme filet de sécurité manuel} : si la
chaîne d'une phase s'interrompt à un point donné (échec de convergence, panne), ce script
permet de relancer \emph{uniquement ce point} à la main (en récupérant
\texttt{CONTCAR}/\texttt{WAVECAR}/\texttt{CHGCAR} du point voisin déjà terminé) sans redemander
un nœud pour rejouer toute la série depuis le début. Ce n'est pas le chemin utilisé par
\texttt{submit\_chain.sh}.

\texttt{job.sh} d'un point (exemple, R-3m à 40~GPa — repli manuel) :
\lstinputlisting{../../hp_curves/Al4C3_R-3m/P40/job.sh}

La soumission des \textbf{7 jobs} (un \texttt{run\_chain.sh} par phase) est orchestrée par
\texttt{hp\_curves/submit\_chain.sh} :
\lstinputlisting{../../hp_curves/submit_chain.sh}

\subsection{Post-traitement : enthalpie et normalisation}

VASP imprime directement l'enthalpie sous contrainte externe dans l'\texttt{OUTCAR} lorsque
\texttt{PSTRESS} $\neq 0$ (ligne \texttt{"enthalpy is~~TOTEN~=~... eV~~P~V=~..."}), qui
correspond exactement à $H = E + PV$ — c'est cette valeur qui est directement exploitée, sans
recalcul. Au point \textbf{0~GPa} (\texttt{PSTRESS=0}), VASP n'imprime pas cette ligne
(puisque $PV=0$, $H=E$) : \texttt{extract\_hp.py} retombe alors sur la ligne standard
\texttt{"free~~energy~~~TOTEN~=~..."} de l'\texttt{OUTCAR}, équivalente à $H$ dans ce cas
particulier. Comme les sept polymorphes n'ont pas le même nombre de formules Al$_4$C$_3$ par
maille (2, 3 ou 4), l'enthalpie est normalisée par formule unitaire (7 atomes) avant
comparaison entre phases. Le script \texttt{hp\_curves/extract\_hp.py} parcourt les 77
répertoires, extrait $H$ et le nombre d'atomes, écrit \texttt{results\_hp.csv} (toutes les
lignes, avec un statut \texttt{converged / unconverged / pending}) ainsi qu'un fichier
\texttt{.dat} par phase (uniquement les points convergés) prêt à être tracé avec
\texttt{pgfplots} (\texttt{\textbackslash addplot table}).

À ce stade (aucun calcul soumis), \texttt{extract\_hp.py} rapporte $0/77$ points convergés —
c'est l'état attendu avant lancement.

\section{Points d'attention}

\begin{itemize}
  \item \textbf{Double détermination du point 0~GPa} : la campagne précédente avait déjà calculé
  un point à 0~GPa, mais à \texttt{KSPACING=0.03} (plus fin) et sans chaînage WAVECAR/CHGCAR
  depuis 50~GPa. Cette campagne calcule un \emph{second} point 0~GPa, cohérent avec le reste de
  la série (\texttt{KSPACING=0.2}, chaîné depuis 10~GPa). Les deux résultats à 0~GPa pourront
  être comparés comme vérification croisée (sensibilité au maillage \textbf{k}) ; c'est le point
  \texttt{KSPACING=0.2} de cette campagne qui sera utilisé pour tracer la courbe H(P) continue,
  par cohérence avec les 10 autres pressions de la même série.
  \item \textbf{\texttt{KPAR=8} vs. le benchmark de parallélisation} : le benchmark réalisé sur
  \texttt{As4S3\_28at} à \texttt{KSPACING=0.2} (27 points \textbf{k} irréductibles, configuration
  la plus proche de celle utilisée ici) avait montré que \texttt{KPAR=4} était \emph{plus
  rapide} que \texttt{KPAR=8} pour ce type de système ($\times1.55$ contre $\times1.13$ par
  rapport à la référence). \texttt{KPAR=8} est néanmoins utilisé ici conformément à la demande
  explicite ; il faut s'attendre à des temps de calcul plus longs qu'avec \texttt{KPAR=4}, sans
  que cela remette en cause la limite de 6~h retenue par point.
  \item \textbf{Chaînage strict, un seul job par phase} : les 11 points d'une phase s'exécutent
  en séquence stricte \emph{à l'intérieur d'un seul job SLURM} — si VASP échoue ou stagne à un
  point donné, \texttt{run\_chain.sh} s'arrête immédiatement (\texttt{exit 1}) et libère le
  nœud sans lancer les points suivants de la phase (les 6 autres phases, chacune sur son propre
  nœud, ne sont pas affectées). Contrairement à la campagne précédente (jobs indépendants
  chaînés par \texttt{--dependency=afterok}), il n'y a plus de risque de queue SLURM entre deux
  points, mais aussi plus de reprise automatique : une reprise manuelle avec le \texttt{job.sh}
  du point interrompu (à l'image de ce qui a été fait pour \texttt{Al4C3\_Pnma\_0GPa} et
  \texttt{Al4C3\_anti-CaFe2O4\_0GPa} dans la campagne précédente) sera nécessaire pour terminer
  la branche, puis relancer manuellement les points suivants un par un ou reprendre
  \texttt{run\_chain.sh} à partir du point qui a échoué.
  \item \textbf{Restart électronique à maillage variable} : si le nombre de points \textbf{k}
  irréductibles change légèrement d'un pas de pression à l'autre (la maille se contracte avec
  la pression), VASP régénère automatiquement les fonctions d'onde/densité de charge manquantes
  à partir du \texttt{WAVECAR}/\texttt{CHGCAR} disponibles (comportement standard
  d'\texttt{ISTART=1}/\texttt{ICHARG=1} en cas de maillage incompatible) ; ce n'est pas bloquant
  mais peut légèrement réduire le bénéfice du redémarrage à certains pas.
\end{itemize}

\section{État d'avancement}

\begin{itemize}
  \item[\textbf{Fait}] Arborescence des 77 répertoires générée
  (\texttt{hp\_curves/gen\_hp\_runs.py}), avec \texttt{POSCAR}/\texttt{POTCAR} des 7 points
  graines copiés depuis les relaxations 50~GPa existantes, et \texttt{INCAR.relax}/\texttt{job.sh}
  écrits pour les 77 points (chaînage \texttt{WAVECAR}+\texttt{CHGCAR}+\texttt{CONTCAR} inclus).
  \item[\textbf{Fait}] Script d'orchestration des soumissions avec dépendances SLURM
  (\texttt{hp\_curves/submit\_chain.sh}) et script d'extraction/normalisation des résultats
  (\texttt{hp\_curves/extract\_hp.py}), tous deux testés (extraction : $0/77$ convergés, comme
  attendu avant tout lancement).
  \item[\textbf{En attente}] Soumission effective des 77 jobs sur le cluster —
  \textbf{en attente de votre validation}.
  \item[\textbf{À venir}] Une fois la campagne terminée : ré-exécution de
  \texttt{extract\_hp.py}, puis tracé de la figure $H(P)$ pour les sept phases entre 0 et
  100~GPa et détermination du domaine de stabilité de chaque polymorphe (enveloppe inférieure
  des courbes $H(P)$, i.e. la phase de plus basse enthalpie à chaque pression), dans un rapport
  de résultats séparé.
\end{itemize}

\end{document}
